\chapter{Embedded Code}\label{sec:EmbeddedCode}
Sometimes, it is necessary to embed code from another (programming) language into the \Java source code\footnote{If you write a code generator tool, this could even by Java code.}.

First, only relativly small fragments of foreign code should be embedded into the Java source. I consider something around 15 to 20~lines still as \bsq{small} in this context, but this can be discussed for sure.

Not only that the embedded code may clutter the Java source, it is also much better modifiable when stored in a resource file\index{resource file}. Therefore larger fragments and full documents can be handled better when provided as resources\index{resource file} (refer to \tqfullvref{sec:ResourcesClasspath})

\keeplisting{7}
Embedded code should be placed into a text block\index{text block}\footnote{For the definition of a text block\index{text block} in Java see \autocite{ORACLE_DOC_LANGUAGE_SPECIFICATION:TextBlocks}.}:
\begin{lstlisting}
final var code = """
````<line1>
````    <line2>
````        <line3>
````    <line4>
````""";
\end{lstlisting}

The left margin of the text inside the text block\index{text block} is indicated by the position of the closing three double quotes. In the sample above, \texttt{<line1>} in line~2 is not indented at all, while the lines~3 and 5 are indented by four spaces, and \texttt{<line3>} is indented by eight spaces. The indentation\index{indentation} within the source code does not matter here; it is just important for optical reasons. The \lstinline|" "| indicates the indentation blanks for the string only.

\keeplisting{5}
The embedded code looks like this when printed:
\begin{lstlisting}[language={}]
<line1>
    <line2>
        <line3>
    <line4>
\end{lstlisting}

Aligning the opening and the closing delimiters (the three double quotes) is not recommended. This would require reindentation of the text block\index{text block} if the variable name or modifiers are changed.

According to chapter \tqfullref{sec:IndentationCharacter} only blanks are used for the indentation\index{indentation!character} of the Java source code. But it is of course possible that the embedded code requires tabs for its indendation.
\keeplisting{6}
In this case, you have to the use the respective escape sequence \bdq{\lstinline|\\t|}:
\begin{lstlisting}[numbers=left]
final var codeJavaScript = """
````\tsetTimeout(function() {
````\t\tmyDisplayer("Hello!");
````\t}, 3000);
````""";
\end{lstlisting}

\bdq{Programmer's Guide to Text Blocks} \autocite{Laskey:TextBlocks} is a comprehensive guide to text blocks\index{text block} in \Java. Text blocks will also be discussed more generally again in chapter \tqfullvref{sec:TextBlocks}.

\subsubsection{Principles}
\begin{enumerate}[label=P\arabic*.]
\item{The opening three double-quotes (the \bsq{\textit{delimiter}}) for a text block\index{text block} should be placed at the right end of the previous line, and the closing delimiter should be placed on its own line, at the left margin of the text block.

Use the line-terminator escape sequence\footnote{A backslash immediately followed by a linefeed.} if the trailing line-feed is not wanted.}
\item{Avoid aligning the opening and closing delimiters and the text block's left margin.}
\item{If tabs are needed for the indentation of the code inside the text block, use the escape sequence \bdq{\lstinline|\\t|}. 

Mixing white space will lead to a result with irregular indentation.}
\item{Do not embed large code blocks.

Use resources for extensive code.}
\end{enumerate}

%------------------------------------------------------------------------------

\section{Formatting Embedding Code}\label{sec:FormattingEmbeddingCode}
If you have already defined formatting rules for the respective type of code you need to embed into your Java source code, you should first check if they can be applied to the embedded code, or if they can be modified to fit for this purpose. If so, use these rules and forget the following sections in this chapter.

If you have other types of code that need to be embedded on a regular basis, you should amend this document accordingly\footnote{Refer to \tqfullref{sec:HowToUseThisDocument}.}, with at least some basic rule on how to format the embedded code of that type.

%------------------------------------------------------------------------------

\subsection{Formatting SQL inside Java}\label{sec:FormattingSQLInsideJava}
If your code needs to access a database\index{database} on an RDBMS\index{RDBMS}, it is quite likely that you have to write some SQL\index{SQL} \autocite{ANSI:SQLStandard,W3SCHOOLS:SQL} statements into your code. Even when you are using an ORM\index{ORM} like JPA\index{JPA} \autocite{ECLIPSE_JPA} and Hibernate\index{Hibernate} \autocite{HIBERNATE_ORM}, it might still be necessary to add some SQL\index{SQL} statements or statements in the tool's variant of SQL (HQL\index{HQL}\index{HQL|see {Hibernate}} in case of Hibernate\index{Hibernate}).

Also the statements of several other query languages can be formatted in the same manner than SQL\index{SQL} statements.

\keeplisting{6}
The SQL \verb#SELECT#\index{SQL!SELECT} statement
\begin{lstlisting}[language=SQL]
SELECT person.id AS "id", person.name AS "lastName", person.surname AS "firstName", address.street AS "street", address.city AS "city"
    FROM contact.person, contact.address
    WHERE person.addressid = address.id
        AND person.profession like '%Teacher%';
\end{lstlisting}

\keeplisting{8}
would be embedded into a Java source file like this:
\begin{lstlisting}
final var sqlSelect = """
````SELECT person.id AS "id", person.name AS "lastName", person.surname AS "firstName", address.street AS "street", address.city AS "city"
````    FROM contact.person, contact.address
````    WHERE person.addressid = address.id
````        AND person.profession like '%Teacher%';\
````""";
\end{lstlisting}

Basically, the keywords are written all uppercase (\verb#SELECT#\index{SQL!SELECT}, \verb#FROM#\index{SQL!FROM},~…), while the arguments are all lowercase (\verb#person.name#, \verb#contact.address#,~…). Alias\index{SQL!alias} names are \textit{always} surrounded by double quotes even they do not contain blanks or are SQL\index{SQL} keywords\footnote{Some RDBMS\index{RDBMS} products do not allow double quotes for this purpose, they require square brackets.}.

The indendation\index{indentation} (four spaces)\index{indendation!character}\index{indentation!size} is per level: \verb#FROM#\index{SQL!FROM} and \verb#WHERE#\index{SQL!WHERE} are one level down from \verb#SELECT#\index{SQL!SELECT}, \verb#AND#\index{SQL!AND} is one level down from \verb#WHERE#\index{SQL!WHERE}.

%==============================================================================
% End of file!!
%==============================================================================

\subsection{Formatting HTML inside Java}\label{sec:FormattingHTMLInsideJava}
Embedded HTML\index{HTML} snippets are usually used to format the output of a program or an application, not only in the context of web applications.

Below some recommendations that were translated from \autocite{selfHTML:GuterHTMLStil}:
\begin{itemize}
	\item{Your HTML markup should be error-free (valid) so that it can be read by browsers and other parsers, such as screen readers. HTML\index{HTML} can only be validated if you include a \verb#doctype# declaration\footnote{This is only relevant if you embedd a complete HTML document into your Java code, and not just HTML snippets.}.} 
	\item{Virtues borrowed from XML\index{XML}: 
		\begin{itemize}
			\item{Well-formed syntax improves readability and enables processing as XHTML\index{XHTML}.}
			\item{All element names, attribute names, and their values are written in lowercase. Attribute values are always enclosed in doubel quotation marks.}
			\item{Use \bdq{meaningful} class names that describe function rather than visual styling.} 
			\item{Elements that do not strictly require a closing tag in HTML (e.g., \verb#p#, \verb#th#, \verb#td#, \verb#dt#, \verb#dd#, \verb#li#) should always be written with a closing tag.}
			\item{Use indentation and blank lines to structure your code clearly.}
		\end{itemize}}
\end{itemize}
\textbf{Note:} Since XHTML\index{XHTML} is an XML\index{XML} dialect, the points listed above are mandatory in XHTML\index{XHTML}, whereas in HTML\index{HTML} they are \bdq{merely} a matter of good practice.

\keeplisting{19}
This short HTML\index{HTML} document
\begin{lstlisting}[language=HTML]
<!DOCTYPE html>
<html>
    <head>
        <meta charset="UTF-8" />
        <title>Title</title>
    </head>
    <body>
        <h1>Headline</h2>
        <p>Text …</p>
        <p>More text …</p>
        <ol>
            <li>Entry 1</li>
            <li>Entry 2</li>
            <li>Entry 3</li>
        </ol>
        <p>Final text!</p>
    </body>
</html>
\end{lstlisting}

\keeplisting{21}
would be embedded like this:
\begin{lstlisting}
final var htmlDocumentString = """
````<!DOCTYPE html>
````<html>
````    <head>
````        <meta charset="UTF-8" />
````        <title>Title</title>
````    </head>
````    <body>
````        <h1>Headline</h2>
````        <p>Text …</p>
````        <p>More text …</p>
````        <ol>
````            <li>Entry 1</li>
````            <li>Entry 2</li>
````            <li>Entry 3</li>
````        </ol>
````        <p>Final text!</p>
````    </body>
````</html>
````""";
\end{lstlisting}

\keeplisting{11}
If you have to embed an HTML fragment – most probably, this would be a \verb#<p># element – this would look like this:
\begin{lstlisting}
final var htmlFragment = """
````<p>Text …</p>
````<p>More text …</p>
````    <ol>
````        <li>Entry 1</li>
````        <li>Entry 2</li>
````        <li>Entry 3</li>
````    </ol>
````<p>Final text!</p>
````""";
\end{lstlisting}

\keeplisting{4}
Or for an \verb#<li># element:
\begin{lstlisting}
final var htmlFragment = """
````<li>Entry 1</li>
````""";
\end{lstlisting}

This means that the top level element will not be indented. If you need to compose an HTML\index{HTML} document from the snippets, jyou can correct the indentation through calling \lstinline|String::indent|\index{java.lang!String!indent()} \autocite{ORACLE_DOC_STRING:indent}. This is explained in more detail in chapter \tqfullvref{sec:TextBlocks}.

%==============================================================================
% End of file!!
%==============================================================================


\subsection{Formatting XML inside Java}\label{sec:FormattingXMLInsideJava}
These formatting rules are valid for plain XML\index{XML}, as well as for various XML\index{XML} dialects, like SVG\index{SVG}, DocBook\index{DocBook} or XSLT\index{XSLT}, but not for XHTML\index{XHTML} – that is covered by \tqfullref{sec:FormattingHTMLInsideJava}, above.

The formatting of XML\index{XML} is relatively straight forward: the top level element will not be indented, each enclosed element will be indented by one level (or four spaces). The first attribute follows immediately the tag, further attributes are aligned with the first one. Attribute values are enclosed in double quotes. Closing tags are aligned with the opening tags, and the closing \bdq{\lstinline| />|} for an empty element is on the same line as the last attribute (or the opening tag, if there are no attributes).

\keeplisting{12}
This would look like this:
\begin{lstlisting}[language=XML]
<element attribute = "value">
    <inner attribute1 = "value"
           attribute2 = "value"
        <empty />
        <other attribute1 = "value"
               attribute2 = "value" />
    </inner>
</element>
<element attribute = "value">
    <inner attribute1 = "value" />
</element>
\end{lstlisting}

\keeplisting{14}
Embedded into Java code, the result would be this:
\begin{lstlisting}
final var xmlSnippet = """
````<element attribute = "value">
````    <inner attribute1 = "value"
````           attribute2 = "value"
````        <empty />
````        <other attribute1 = "value"
````               attribute2 = "value" />
````    </inner>
````</element>
````<element attribute = "value">
````    <inner attribute1 = "value" />
````</element>
````""";
\end{lstlisting}

%==============================================================================
% End of file!!
%==============================================================================


%==============================================================================
% End of file!!
%==============================================================================


%==============================================================================
% End of file!!
%==============================================================================
